% ──────────────────────────────────────────────────────────────────────────────
% DISCUSSION SECTION DRAFT
% ──────────────────────────────────────────────────────────────────────────────

\section{Discussion}
\label{sec:discussion}

\subsection{Interpretation of Main Findings}

The central result — that reweighting-based methods degrade sharply under
positive-class scarcity while the proposed framework does not — has a
straightforward structural explanation.
Group DRO and OT Repair both operate exclusively on the existing training
distribution: Group DRO up-weights under-represented groups; OT Repair
redistributes mass between group-conditional distributions.
When the disadvantaged group retains only 43 positive examples ($p_+ = 10\%$),
there is insufficient signal for either method to construct meaningful
group-level gradients or a reliable transport plan.
The RL framework sidesteps this constraint entirely by generating new minority
positive samples rather than reweighting existing ones — making it the only
approach that can improve representation in the region where the classifier
currently fails.

The credit card dataset is included to assess generalisability beyond the
primary fairness setting.
Because the no-intervention baseline already achieves near-zero EO on this
dataset, the method has no meaningful fairness gap to close, and the results
reflect this: EO changes across all methods are small and within the margin
of uncertainty.
On utility, the RL framework records the lowest AUC at both bias levels,
suggesting a marginal cost from generating synthetic samples when the
classifier has no clear decision boundary to be guided toward.
This is consistent with the framework's design — it is optimised for the
positive-class scarcity regime, and its benefit is proportional to the
size of the gap that exists.
The credit results therefore serve to delineate the applicability boundary:
the framework is most valuable when a genuine fairness gap is present and
should be expected to offer limited or no benefit in settings where the
baseline classifier is already near-fair.

CT-GAN's high variance at $p_+ = 10\%$ reflects a known weakness of
conditional generative models at very low sample counts.
With only 43 minority positive examples available to condition on, CT-GAN's
training is unstable across seeds, leading to inconsistent generation quality.
The RL framework, by operating in PCA-compressed space with bounded
perturbations from real seed points, is less sensitive to this regime because
it does not need to learn a full generative model from scratch — it perturbs
existing examples rather than generating unconditionally.

\subsection{Effect of Design Choices}

The ablation results collectively identify three design choices as load-bearing
for the framework's effectiveness.

\textbf{PCA compression.}
Reducing the action space from the full feature dimensionality (${\sim}100$)
to 10 PCA components is the single most impactful design choice, producing a
$5\times$ reduction in EO gap.
In high-dimensional spaces, bounded delta actions cover an exponentially small
fraction of the neighbourhood around each seed point, making directed
exploration of the decision boundary intractable.
PCA concentrates the principal axes of variation in the minority group's
training distribution into a low-dimensional subspace, so each delta action
moves the generated sample along a semantically meaningful direction rather
than an arbitrary raw-feature direction.

\textbf{Credit assignment ($\gamma = 1.0$).}
The switch from $\gamma = 0.99$ to $\gamma = 1.0$ was the single change that
eliminated the training deadzone — reducing the fraction of episodes where
$\beta$ underperforms $\alpha$ from roughly 50\% to 10--12\%.
With $\gamma = 0.99$ and a trajectory length of $T = 2000$, the effective
credit horizon is $1/(1-\gamma) = 100$ steps; steps 500--2000 contribute a
discount factor below $0.007$ and are effectively ignored in the policy
gradient.
Since the beta classifier is trained on the full 2000-step trajectory, this
creates a systematic mismatch: the reward signal used to update the policy does
not reflect the quality of the majority of the generated samples.
Setting $\gamma = 1.0$ removes the discount entirely, so all 2000 trajectory
steps contribute equally to the policy gradient and the deadzone resolves.

\textbf{Delta action scale.}
The $4\times$ EO degradation at both $\delta = 0.05$ and $\delta = 0.20$
relative to the default $\delta = 0.10$ reflects a narrow operating regime.
Too small a perturbation places generated samples near the training data
manifold, where the beta classifier is already confident and the DVRL local
reward signal is near zero.
Too large a perturbation overshoots the current decision boundary, placing
samples in out-of-distribution regions that do not provide useful gradient
information and can actively destabilise beta's training.
The optimal $\delta = 0.10$ keeps generated samples within the classifier's
uncertainty region — consistently near the boundary but not beyond it.

\subsection{Strengths}

The framework's primary strength is its applicability to a regime that existing
methods cannot address.
It achieves the best fairness-utility tradeoff on the primary dataset under
high scarcity, matching or exceeding all baselines on both EO and utility
simultaneously.
Unlike reweighting methods, it does not require a sufficient number of
minority positive examples to function — it can improve a classifier even when
only a handful of such examples exist in the training set.

The framework also shows no degradation on unbiased data.
Applying it to a training set where the classifier is already near-optimal does
not harm predictive utility, which is a necessary condition for safe deployment
in practice — a method that only works when data are biased and degrades
otherwise is of limited practical value.

The PCA-plus-delta action design is also computationally efficient relative to
full generative models.
CT-GAN and similar approaches require training a separate generative model,
which is expensive and unstable at low sample counts.
The RL framework needs only a PCA fit and bounded perturbations, making it
deployable in low-data regimes where full generative models are impractical.

\subsection{Limitations}

\textbf{Seed variance.}
The most significant practical limitation is the persistent per-seed variance.
Across five seeds, a minority of runs degrades EO relative to the no-intervention
baseline.
The rogue-seed pattern is most pronounced at $p_+ = 5\%$, where only 17
minority positive examples remain.
At this extreme level of scarcity, the agent's ability to find the correct
decision boundary region is sensitive to the initialisation of the beta
classifier, and the small number of real positive examples limits the diversity
of starting points for delta action trajectories.
The mean results are robust, but the seed range is wide enough that individual
deployments without ensemble averaging may not reliably improve fairness.

\textbf{Low-baseline-EO datasets.}
As demonstrated on credit card, the framework provides no benefit when the
baseline classifier already achieves near-zero fairness gap.
It is specifically designed for the positive-class scarcity regime, and its
contribution is proportional to the size of the gap it has to close.
Practitioners should assess the baseline EO gap before applying the framework.

\textbf{Hyperparameter sensitivity.}
The delta action scale $\delta$ exhibits a narrow operating range where
performance is optimal, and this range is likely dataset-dependent.
The current experiments use a fixed $\delta = 0.10$ tuned on Census; transfer
to new datasets may require re-tuning, and the sensitivity to this parameter
adds an additional calibration step relative to methods that have fewer
hyperparameters.

\textbf{Computational cost.}
The framework retrains the beta classifier from scratch at each episode across
$T = 800$ training episodes.
Although each episode is lightweight, the cumulative cost scales linearly with
the number of episodes and the classifier training depth.
On the datasets used here this is manageable, but scaling to higher-dimensional
feature spaces or larger training sets would require either a faster beta
classifier or a reduced episode count.

\end{document}
